\documentclass[11pt,a4paper]{article}
\usepackage{brainteaser}

\begin{document}
\brainteasertitle{7}{La fausse pièce}

\begin{statementbox}
On possède $12$ pièces visuellement identiques. Une seule est fausse, mais on
ignore laquelle et si elle est plus lourde ou plus légère que les autres.

Une balance à deux plateaux donne trois résultats: gauche plus lourde,
équilibre ou droite plus lourde. Identifier avec certitude la fausse pièce et
déterminer si elle est plus lourde ou plus légère en seulement trois pesées.
\end{statementbox}

\begin{hintsbox}
\begin{enumerate}
  \item Combien de situations faut-il distinguer? Combien de feuilles un
  arbre de trois pesées peut-il avoir au maximum?
  \item Pourquoi une première pesée de trois pièces contre trois crée-t-elle
  une branche difficile?
  \item Commencer plutôt par peser quatre pièces contre quatre.
  \item Si cette pesée est équilibrée, quelles pièces sont innocentes?
  \item Si elle penche à gauche, quelles sont exactement les huit possibilités
  encore compatibles avec ce résultat?
\end{enumerate}
\end{hintsbox}

\correction

Il existe $12\times2=24$ situations possibles. Trois pesées donnent au plus
$3^3=27$ suites de résultats: une solution est théoriquement possible.
Numérotons les pièces de $1$ à $12$.

\subsection*{Première pesée}
Peser
\[
1,2,3,4\quad\text{contre}\quad5,6,7,8.
\]

\subsection*{Cas 1: équilibre}
La fausse pièce est parmi $9,10,11,12$, et les pièces $1$ à $8$ sont vraies.
Peser
\[
9,10,11\quad\text{contre}\quad1,2,3.
\]
\begin{itemize}
  \item Équilibre: la pièce $12$ est fausse. La peser contre une vraie pièce
  indique si elle est lourde ou légère.
  \item Gauche plus lourde: une des pièces $9,10,11$ est trop lourde. Peser
  $9$ contre $10$.
  \item Gauche plus légère: une des pièces $9,10,11$ est trop légère. Peser
  $9$ contre $10$.
\end{itemize}
Dans les deux derniers cas, le déséquilibre désigne la fausse pièce et un
équilibre désigne la pièce $11$.

\subsection*{Cas 2: le premier plateau est plus lourd}
Les possibilités sont
\[
1_H,2_H,3_H,4_H,5_L,6_L,7_L,8_L,
\]
où $H$ signifie «trop lourde» et $L$ «trop légère». Les pièces $9$ à $12$ sont
authentiques. Peser ensuite
\[
1,2,5\quad\text{contre}\quad3,6,9.
\]
\begin{itemize}
  \item \textbf{Gauche plus lourde:} possibilités $1_H,2_H,6_L$. Peser $1$
  contre $2$. Le plateau plus lourd désigne la fausse; l'équilibre désigne
  $6_L$.
  \item \textbf{Droite plus lourde:} possibilités $3_H,5_L$. Peser $3$
  contre la vraie pièce $9$. Si $3$ est plus lourde, c'est $3_H$; sinon,
  c'est $5_L$.
  \item \textbf{Équilibre:} possibilités $4_H,7_L,8_L$. Peser $7$ contre
  $8$. La plus légère est fausse; l'équilibre désigne $4_H$.
\end{itemize}

Si la première pesée penche dans l'autre sens, appliquer le raisonnement
symétrique en échangeant «lourde» et «légère».

\begin{resultbox}
Cet arbre de décision identifie toujours la pièce et la nature de l'anomalie
en trois pesées.
\end{resultbox}

\end{document}
